\documentclass[border=3pt]{standalone}
% 公共导言区 —— 所有 TikZ 图 \input 本文件后使用
\usepackage[UTF8, fontset=mac]{ctex}
\usepackage{tikz}
\usetikzlibrary{arrows.meta, positioning, calc, shapes.geometric, decorations.pathmorphing, backgrounds, fit}
\definecolor{zhusha}{HTML}{9B2D20}
\definecolor{mohei}{HTML}{1A1A1A}
\definecolor{shenhui}{HTML}{555555}
\definecolor{qianhui}{HTML}{F5F0EC}
\definecolor{lan}{HTML}{1F4E79}
\definecolor{qing}{HTML}{2E7D6E}
\tikzset{
  block/.style={draw=shenhui, fill=white, thick, rounded corners=1.5pt,
                font=\small, align=center, inner sep=4pt, minimum height=7mm},
  core/.style={block, fill=lan!12, draw=lan, font=\small\bfseries},
  periph/.style={block, fill=qing!10, draw=qing},
  power/.style={block, fill=zhusha!8, draw=zhusha},
  note/.style={font=\footnotesize, text=shenhui, align=center},
  lab/.style={font=\footnotesize, text=mohei},
  sig/.style={-{Stealth[length=2.2mm]}, thick, color=mohei},
  fbsig/.style={-{Stealth[length=2.2mm]}, thick, color=zhusha, dashed},
  vec/.style={-{Stealth[length=3mm]}, very thick, color=zhusha},
}

\begin{document}
\begin{tikzpicture}[x=1mm, y=1mm, font=\small]

% ===== 母线 =====
\draw[very thick] (-19, 34) -- (19, 34);
\node[lab, anchor=south] at (0, 35.5) {$V_{dc}$ 母线（+）};
\draw[very thick] (-19, 2) -- (19, 2);
\node[lab, anchor=north] at (0, 0.5) {母线（$-$）};

% ===== 三相桥腿（Q+ / 中点 / Q- / 采样电阻）=====
\foreach \x/\s in {-18/1, 0/2, 18/3} {
  \draw[thick, fill=white] (\x-3.2, 25) rectangle (\x+3.2, 31);
  \node[note] at (\x, 28) {$Q_{\s}^{+}$};
  \draw[thick] (\x, 34) -- (\x, 31);
  \draw[thick] (\x, 25) -- (\x, 13);          % 中点段
  \draw[thick, fill=white] (\x-3.2, 7) rectangle (\x+3.2, 13);
  \node[note] at (\x, 10) {$Q_{\s}^{-}$};
  \draw[thick] (\x, 7) -- (\x, 4.5);
  \draw[thick, fill=zhusha!20, draw=zhusha] (\x-2.5, 2.9) rectangle (\x+2.5, 4.5);
}

% ===== 相线：从桥臂中点侧向引出，绕开负母线两端后下行 =====
% U 相：左引出，x=-24（母线左端 -19 之外）
\draw[thick, lan] (-18, 19) -- (-24, 19) -- (-24, -5);
\node[fill=mohei, circle, inner sep=0.9pt] at (-18, 19) {};
% V 相：同样左引，更低一层走 x=-32
\draw[thick, lan] (0, 16) -- (-32, 16) -- (-32, -9);
\node[fill=mohei, circle, inner sep=0.9pt] at (0, 16) {};
% W 相：右引出，x=24（母线右端 19 之外）
\draw[thick, lan] (18, 19) -- (24, 19) -- (24, -5);
\node[fill=mohei, circle, inner sep=0.9pt] at (18, 19) {};
\node[note, anchor=south east, text=lan] at (-24.8, 20) {$U$};
\node[note, anchor=south east, text=lan] at (-32.8, 17) {$V$};
\node[note, anchor=south west, text=lan] at (24.8, 20) {$W$};

% ===== 三相绕组（星接，圆圈内标相名）=====
\foreach \x/\y/\n in {-24/-5/U, -32/-9/V, 24/-5/W} {
  \draw[very thick, lan, fill=lan!8] (\x, \y) circle (3.2);
  \node[lab, text=lan, font=\small\bfseries] at (\x, \y) {\n};
}
% 中性点
\coordinate (star) at (-8, -14);
\draw[thick] (-24, -8.2) -- (star);
\draw[thick] (-32, -12.2) -- (star);
\draw[thick] (24, -8.2) -- (star);
\node[circle, fill=mohei, inner sep=1.1pt] at (star) {};
\node[lab, anchor=north] at (-8, -15.5) {中性点 $O$};

% ===== 注记 =====
\node[note, anchor=east, align=right] at (-30, 3.7) {低侧采样电阻\\（三电阻检测）};
% 母线电容（右侧，普通导线）
\draw[thick] (19, 34) -- (38, 34);
\draw[thick] (38, 34) -- (38, 21);
\draw[thick] (34, 21) -- (42, 21);
\draw[thick] (34, 17) -- (42, 17);
\draw[thick] (38, 17) -- (38, 2);
\draw[thick] (19, 2) -- (38, 2);
\node[note, anchor=west] at (43, 19) {母线电容};

\end{tikzpicture}
\end{document}
